%
% Capítulo 8
%
\chapter{Estado Atual da Beta e Próximos Passos} \label{cap:estadobeta}

A versão beta do projeto representa uma fase intermédia de consolidação funcional da solução,
em que os principais fluxos do sistema já se encontram implementados e testados, embora ainda existam aspetos a melhorar e integrar totalmente.

A solução já não se encontra limitada a um protótipo exclusivamente local. Com a introdução do backend em Ktor, a integração com PostgreSQL e a utilização de uma API REST funcional, o sistema passou a assumir uma estrutura mais próxima da versão final pretendida.

\section{Funcionalidades Concluídas na Versão Beta}

Na fase atual do projeto, já se encontram implementadas as funcionalidades centrais do sistema.
Entre os principais elementos já concluídos destacam-se:
\begin{itemize}
    \item aplicação Android com separação de fluxos entre professor e aluno;
    \item gestão de sessão e adaptação da interface ao perfil autenticado;
    \item registo e autenticação de professores e alunos;
    \item associação, desassociação e troca de professor por parte do aluno;
    \item definição e persistência de disponibilidades, incluindo múltiplos intervalos no mesmo dia;
    \item definição e persistência de restrições do professor;
    \item criação de propostas de horário;
    \item visualização organizada do horário semanal criado;
    \item integração do horário aceite com o Google Calendar do professor;
    \item backend funcional com API REST;
    \item persistência remota em PostgreSQL;
    \item validação dos endpoints principais com Postman.
\end{itemize}

Estes resultados demonstram que a versão beta já suporta os casos de uso mais importantes do projeto, permitindo recolher a informação necessária, processá-la e apresentar ao utilizador uma proposta concreta de horário.

\section{Estado da Integração do Sistema}

Um dos principais avanços da versão beta foi a integração progressiva entre frontend e backend.
A aplicação Android já comunica com a API para um conjunto importante de operações, nomeadamente na gestão de utilizadores,
autenticação, disponibilidades, restrições e criação de horários.
Esta evolução permitiu validar a arquitetura distribuída da solução e reduzir a dependência de mecanismos exclusivamente locais.

Apesar deste progresso, a integração ainda não pode ser considerada totalmente concluída em todos os fluxos.
Em alguns pontos, subsistem aspetos híbridos entre persistência local e remota, o que exige atenção adicional para garantir consistência total
entre os diferentes componentes. Ainda assim, a base da comunicação entre cliente e servidor encontra-se estabelecida e operacional,
constituindo um avanço significativo face às fases iniciais do projeto.

\section{Limitações Atuais}

Embora a versão beta demonstre a viabilidade da solução, existem ainda limitações que deverão ser resolvidas até à entrega final.
Uma dessas limitações prende-se com a necessidade de consolidar plenamente a sincronização entre os dados locais armazenados em Room e os dados persistidos no backend.
A coexistência destas duas camadas de persistência, embora útil durante o desenvolvimento incremental, introduz complexidade adicional e requer uma definição mais clara do papel de cada uma na arquitetura final.

Outra limitação atual diz respeito ao algoritmo de criação de horários.
A abordagem implementada já é suficiente para validar os fluxos principais e construir propostas coerentes em cenários representativos,
mas poderá ainda ser refinada para responder melhor a casos mais complexos, a preferências adicionais ou a critérios mais exigentes de distribuição das sessões.

Para além disso, existem aspetos a melhorar ao nível do polimento da interface, da robustez de alguns fluxos e da cobertura de testes.
Mantém-se também como linha possível de evolução a avaliação de mecanismos adicionais de entrada de dados, nomeadamente com apoio de modelos multimodais para interpretação de texto, imagem e áudio, conforme descrito no capítulo dedicado a esse estudo.

\section{Próximos Passos}

Tendo em conta o estado atual do projeto, os próximos passos mais relevantes passam por consolidar a solução já implementada e preparar a transição para a versão final.
Entre as prioridades identificadas destacam-se:
\begin{itemize}
    \item consolidar a integração entre frontend e backend nos fluxos ainda híbridos;
    \item reduzir as dependências entre persistência local e remota;
    \item reforçar a coerência e sincronização dos dados do sistema;
    \item melhorar a robustez da criação de horários;
    \item aprofundar os testes funcionais e de integração;
    \item continuar o polimento da interface e da apresentação final da solução;
    \item avaliar a viabilidade de funcionalidades baseadas em modelos multimodais para apoio ao registo de disponibilidades.
\end{itemize}

Adicionalmente, prevê-se a continuação do estudo exploratório de funcionalidades mais avançadas caso tal se revele viável dentro do calendário do projeto.
Esta linha de evolução poderá acrescentar valor significativo à solução, mas deverá ser enquadrada de forma realista face ao estado atual de desenvolvimento.

\section{Balanço da Versão Beta}

De uma forma geral, a versão beta permitiu demonstrar que o projeto já possui uma base funcional consistente e tecnicamente viável. O sistema consegue suportar os fluxos
fundamentais do problema, desde a recolha de dados até à criação e apresentação de propostas de horário. A arquitetura adotada mostrou-se adequada à evolução do projeto e os resultados obtidos nesta fase constituem um indicador positivo para a continuação do trabalho.

A partir deste ponto, o foco deixa de estar apenas na implementação de novas funcionalidades e passa também pela consolidação, integração, validação e refinamento do que já foi desenvolvido.
Este balanço é particularmente importante, pois permite encarar a versão final não como o início da solução, mas como uma etapa de maturação e aperfeiçoamento de uma base já operacional.
